\documentclass[../../main.tex]{subfiles}
\begin{document}

{\bf [解]}

題意の数列が
\begin{align}
  a_n = 2^{n} \label{eq:1}
\end{align}
であることを数学的帰納法により示す．まず，この形は条件(1)を満たすので，以下条件(2)についてだけ考えれば良い．

$n=1$ の時は$a_1=1$より成立する．
そこで$n \leqq k\, (k \in \mathbb{N})$ での成立を仮定すると，
$a_1,\cdots,a_k$は2進数表示で$1,10,\cdots,\underbrace{10\dots 0}_{k}$を表す．
これらを組み合わせると$T = \underbrace{111\dots1}_{k}=1+2+2^{2}+\cdots+2^{k}$ から $1$ までの数を全てつくることができるから，条件(2)より
\begin{align}
  a_{k+1} \ge T + 1 = \sum_{i=1}^{k}2^{i}+1 = 2^{k+1}
\end{align}
である必要がある．題意より$a_{k+1}$は条件(1),(2)を満たす最小の整数だから
\begin{align}
  a_{k+1} = 2^{k+1}
\end{align}
である．よって $n = k+1$ でも\cref{eq:1}は成立するので数学的帰納法により
\begin{align}
  a_n = 2^{n}
\end{align}
である．

{\bf [解説]}

[解]では2進数表示を考えることによって$1, 2, \dots, 2^n$から$1 \sim 2^{n+1}-1$を作れることを示したが，ここにも数学的帰納法を用いても以下のように示せる．

$n=k$ で仮定して$n=k+1$について考える．
$1, 2, \dots, 2^{k+1}$ について
$1 \sim 2^{k+1}-1$ は仮定よりつくれるから，
\begin{align}
  \begin{cases}
    1^\circ & 2^{k+1} \text{ はつくれる} \\
    2^\circ & 2^{k+1}+1, \dots, 2^{k+2}-1 \text{ は } 2^{k+1} \text{ を加えてつくれる}
  \end{cases}
\end{align}
よって $n = k+1$ でも成立して示された．

\end{document}
